\clearpage\begingroup\SourceLeaf{71b}\HistoricalHead{66}{L}{PHYSIK. NACHLASS.}
\begin{center}[8.]\end{center}
\noindent{\small [S. 207]}
\begin{center}\textbf{[Zurückwerfungsmethode.]}\end{center}
\indent Der Lauf der Elongationen, bei abwechselnd angebrachten Inductionsstössen
kann am übersichtlichsten so dargestellt werden:
\par\noindent\hspace*{\parindent}\(\log\dfrac1{\mu\mu}=\text{Decrem. Logarithm.}\)
\par\noindent\hspace*{\parindent}\(x\) Wirkung eines Elongationsstosses
\par\noindent\hspace*{\parindent}Elongationen (zum natürlichen Ruhestand hinzuzusetzen):
\[
\begin{aligned}
&-\frac{\mu^3}{1+\mu^4}x+e\\[1ex]
&-\frac{\mu}{1+\mu^4}x-e\mu\mu\\[1ex]
&+\frac{\mu^3}{1+\mu^4}x+e\mu^4\\[1ex]
&+\frac{\mu}{1+\mu^4}x-e\mu^6\\[1ex]
&-\frac{\mu^3}{1+\mu^4}x+e\mu^8\\[1ex]
&-\frac{\mu}{1+\mu^4}x-e\mu^{10}
\end{aligned}
\]
\begin{center}etc.\end{center}
\clearpage\SourceLeaf{72a}\HistoricalHead{67}{R}{BERECHNUNG MAGNETISCHER KRÄFTE.}
\begin{center}\textbf{BEMERKUNGEN.}\end{center}
\begingroup\small
\indent Der Sinn der vorstehenden Rechnung wird klar, wenn man die Grenze betrachtet, der mit wachsender
Zeit die Ausschläge zustreben. Dann verschwindet offenbar in jeder Gleichung das zweite Glied, und vier
aufeinanderfolgende Werte der Elongationen haben folgende Beträge:
\[
\begin{aligned}
&+\frac{\mu}{1+\mu^4}x,\\[1ex]
&-\frac{\mu^3}{1+\mu^4}x,\\[1ex]
&-\frac{\mu}{1+\mu^4}x,\\[1ex]
&+\frac{\mu^3}{1+\mu^4}x.
\end{aligned}
\]
\noindent Dadurch werden, wie man sieht, auf der Skala vier Punkte bestimmt; der Abstand der beiden äußeren
werde mit \(a\), der der beiden inneren mit \(b\) bezeichnet. \(a\) und \(b\) haben demnach folgende Werte:
\[
\begin{aligned}
a&=\frac{2\mu x}{1+\mu^4},\\[1ex]
b&=\frac{2\mu^3x}{1+\mu^4}.
\end{aligned}
\]
\noindent Das Verhältnis \(\dfrac ab\) wird also gleich \(\dfrac1{\mu^2}\), und man erkennt, daß die hier mitgeteilten \textsc{Gauss}schen Formeln
die Theorie der sogenannten Zurückwerfungsmethode bilden, die \textsc{Gauss} und \textsc{Weber} gemeinschaftlich zugeschrieben
wird, und von \textsc{Weber} (a.\,a.\,O., \textsc{Wilh. Webers} Werke III, S. 441) veröffentlicht worden ist.
Übrigens fehlen in den \textsc{Gauss}schen Formeln höhere Glieder, die in der Tat irrelevant sind, da beim Versuch
der stationäre Zustand immer abgewartet wird.
\par\raggedleft\textsc{Schaefer}.\par
\clearpage\endgroup\endgroup
